% ============================================================================
%  chapter_20_continual_report.tex
%  Chapter 20: The Continual Simulation Report
%  Identity, Scale, and the Kernel Event Spine
%  VSEPR-SIM Theoretical Division  |  v5.13 Branch Freeze Document
%  Compile: pdflatex (twice for TOC / refs)
% ============================================================================
\documentclass[11pt,a4paper]{article}

%% -- Encoding & fonts -------------------------------------------------------
\usepackage[utf8]{inputenc}
\usepackage[T1]{fontenc}
\usepackage{lmodern}
\usepackage{microtype}

%% -- Core packages ----------------------------------------------------------
\usepackage{amsmath,amssymb,amsfonts,mathtools,bm}
\usepackage{geometry}
\usepackage{xcolor}
\usepackage{booktabs,array,tabularx,longtable}
\usepackage{enumitem}
\usepackage{fancyhdr}
\usepackage{titlesec}
\usepackage{hyperref}
\usepackage{float}
\usepackage{caption}
\usepackage{mdframed}
\usepackage{listings}
\usepackage{amsthm}

%% -- Page geometry ----------------------------------------------------------
\geometry{left=2.8cm,right=2.8cm,top=3.2cm,bottom=3.0cm,headheight=15pt}

%% -- Hyperref ---------------------------------------------------------------
\hypersetup{
  colorlinks=true, linkcolor=black, citecolor=black, urlcolor=black,
  pdftitle={Chapter 20 -- The Continual Simulation Report},
  pdfsubject={VSEPR-SIM v5.13 Branch Freeze -- IKK Identity, Scale, Event Spine}
}

%% -- Section numbering: 20.0 through 20.N ----------------------------------
\renewcommand{\thesection}{20.\arabic{section}}
\setcounter{section}{-1}
\renewcommand{\thesubsection}{\Alph{subsection}}

%% -- Section styling --------------------------------------------------------
\titleformat{\section}
  {\normalfont\large\bfseries}{\thesection.}{0.6em}{}
  [\vspace{-2pt}\rule{\textwidth}{0.4pt}]
\titleformat{\subsection}
  {\normalfont\normalsize\bfseries}{\thesubsection.}{0.5em}{}
\titleformat{\subsubsection}
  {\normalfont\normalsize\itshape}{\thesubsubsection.}{0.5em}{}

%% -- Header / footer --------------------------------------------------------
\pagestyle{fancy}
\fancyhf{}
\fancyhead[L]{\small\textit{VSEPR-SIM Chapter 20 --- The Continual Simulation Report}}
\fancyhead[R]{\small\textit{v5.13 Freeze $\cdot$ Day 82}}
\fancyfoot[C]{\thepage}
\renewcommand{\headrulewidth}{0.4pt}

%% -- Colour palette ---------------------------------------------------------
\definecolor{vsimdark}{RGB}{26,82,118}
\definecolor{vsimgreen}{RGB}{30,130,76}
\definecolor{vsimred}{RGB}{192,57,43}
\definecolor{vsimamber}{RGB}{211,84,0}
\definecolor{codegray}{RGB}{245,245,245}

%% -- Theorem environments ---------------------------------------------------
\theoremstyle{definition}
\newtheorem{definition}{Definition}[section]
\newtheorem{axiom}{Axiom}[section]
\newtheorem{principle}{Principle}[section]

%% -- Mdframed boxes ---------------------------------------------------------
\mdfdefinestyle{doctrinebox}{
  linecolor=vsimdark, linewidth=1.2pt,
  backgroundcolor=blue!4,
  innertopmargin=8pt, innerbottommargin=8pt,
  innerleftmargin=10pt, innerrightmargin=10pt
}
\mdfdefinestyle{warnbox}{
  linecolor=vsimamber, linewidth=1.2pt,
  backgroundcolor=yellow!8,
  innertopmargin=8pt, innerbottommargin=8pt,
  innerleftmargin=10pt, innerrightmargin=10pt
}
\mdfdefinestyle{freezebox}{
  linecolor=vsimgreen, linewidth=1.5pt,
  backgroundcolor=green!5,
  innertopmargin=8pt, innerbottommargin=8pt,
  innerleftmargin=10pt, innerrightmargin=10pt
}

%% -- Listings ---------------------------------------------------------------
\lstdefinestyle{eventprint}{
  basicstyle=\ttfamily\small,
  backgroundcolor=\color{codegray},
  breaklines=true, frame=single,
  framesep=4pt, rulecolor=\color{gray!40}
}

%% -- IKK macro set (harmonised with IKK IV v1.0) ---------------------------
\newcommand{\Dfrak}{\mathfrak{D}}
\newcommand{\Ifrak}{\mathfrak{I}}
\newcommand{\etaab}{\eta_{ab}}
\newcommand{\Psihid}{|\Psi^{\mathrm{hid}}\rangle}
\newcommand{\Drec}{D_{\mathrm{rec}}}
\newcommand{\DeltaD}{\Delta\Dfrak}
\newcommand{\DeltaS}{\Delta S}
\newcommand{\Sn}{S_n}
\newcommand{\ivec}{\bar{\imath}_f}

%% -- IKK badge commands ----------------------------------------------------
\newcommand{\Dpass}[1]{\colorbox{green!25}{\texttt{\small #1}}}
\newcommand{\Dwarn}[1]{\colorbox{yellow!45}{\texttt{\small #1}}}
\newcommand{\Dfail}[1]{\colorbox{red!25}{\texttt{\small #1}}}

% ===========================================================================
\begin{document}

\thispagestyle{empty}

\begin{center}
  {\footnotesize\sffamily VSEPR-SIM Theoretical Division
	\quad\textperiodcentered\quad
	\textsc{v5.13 Branch Freeze Document}
	\quad\textperiodcentered\quad Day 82}\\[10pt]
  {\LARGE\bfseries\color{vsimdark} Chapter 20}\\[6pt]
  {\Large\bfseries The Continual Simulation Report}\\[4pt]
  {\large Identity, Scale, and the Kernel Event Spine}\\[16pt]
  {\normalsize\itshape
	IKK End-Tag Enrichment (WO-75A)
	\quad\textperiodcentered\quad
	IKK Identity Vector (WO-75B)
	\quad\textperiodcentered\quad
	v5.13 Freeze Summary}\\[6pt]
  {\small 2026-06-28 \quad|\quad
	Commit: \texttt{489d11cd} \quad|\quad
	167/167 tests passing}
\end{center}

\vspace{12pt}
\begin{mdframed}[style=freezebox]
\textbf{Branch freeze notice.}\quad
This document is the formal closure artefact for the \texttt{v5.13} arc of
VSEPR-SIM\@.  The branch \texttt{feature/wizard-full-module-expansion} is
frozen at this state.  All deferred items are tracked in \texttt{STAGE.md}
and \texttt{CLOSEOUT.md}.  Work resumes on a new branch.
\end{mdframed}

\vspace{8pt}
\tableofcontents
\clearpage

% ===========================================================================
\section{Preamble}
% ===========================================================================

This chapter documents a single organising idea that runs through the entire
v5.13 arc and gives every module within it its architectural justification:
\emph{the simulation must continuously print what it knows}.

The idea is not about logging.  It is not a debugging aid.  It is the
foundational output doctrine of the platform: every event the kernel observes,
every projection loss it incurs, every identity fraction it retains or
discards---all of it must be expressed continuously, per step, in structured,
human-readable, machine-parseable form.

This principle existed implicitly from Day~1 in the design of the
\texttt{KernelEventLog}.  The v5.13 arc formalised it, extended it upward into
the IKK theoretical layer, and delivered two concrete implementations:
the IKK End-Tag system (WO-75A) and the IKK Identity Vector (WO-75B).

% ===========================================================================
\section{The Continual Printing System}
% ===========================================================================

\subsection{Doctrine}

\begin{mdframed}[style=doctrinebox]
\textbf{Core doctrine.}\quad
A simulation that only produces a final result is not a simulation.
It is a calculator.  A simulation must continuously narrate its own state,
uncertainties, and information losses across every step, every scale, and every
event channel---so that the researcher can reconstruct not only \emph{what
happened} but \emph{what the kernel knew}, and \emph{what it could not know},
at every moment of the run.
\end{mdframed}

\vspace{4pt}
This doctrine has three consequences for the architecture:

\begin{enumerate}[leftmargin=2em]
  \item \textbf{No silent discards.}  Every event the kernel processes is
	emitted---to the console, to the JSONL archive, or to both.  The event
	limiter may sample the console under load; the JSONL archive is never
	sampled.

  \item \textbf{Structured streams.}  Each physical channel emits a
	distinguishable, grep-able prefix.  A researcher running a live session can
	follow a specific channel without instrumenting new code.

  \item \textbf{Identity accounting.}  The simulation does not only report
	\emph{what happened}; it reports \emph{how much of what happened is still
	recoverable} in the language of IKK distinguishability
	($\Dfrak$) and cross-scale efficiency ($\etaab$).
\end{enumerate}

\subsection{The Four Live Print Streams}

The runtime emits four structured event channels to the terminal and the
JSONL archive simultaneously.  Each is individually grep-able and forms a
separate column in the JSONL event archive:

\begin{lstlisting}[style=eventprint,
  caption={Representative output from all four live print channels.}]
[PP]   step=184 t=0.184 idA=e_001 idB=pos_001 r=0.018 Erel=0.382 action=ANNIHILATE
[CHEM] step=184 t=0.184 idA=C_004 idB=O_002 r=1.24 type=nonbond action=NONE
[DECAY]step=522 t=0.522 id=Pu_239_001 mode=alpha daughter=U235_001 dE=0.052
[ANN]  step=184 t=0.184 pair=e_001+pos_001 Ein=1.022 Eout=1.022 products=gamma_001,gamma_002
\end{lstlisting}

\begin{table}[H]
\centering
\small
\begin{tabular}{@{}llp{7cm}@{}}
\toprule
Prefix & Channel & Trigger \\
\midrule
\texttt{[PP]}    & Particle--particle interaction
  & Any pair interaction; \texttt{ANNIHILATE} on gate pass \\
\texttt{[CHEM]}  & Chemistry bond / nonbond
  & Bond formation, breaking, or nonbond pair event \\
\texttt{[DECAY]} & Single-particle decay
  & Decay event (alpha, beta, fission) \\
\texttt{[ANN]}   & Anti-pair annihilation
  & Triggered \texttt{AnnihilationEvent} \\
\bottomrule
\end{tabular}
\caption{Live print stream channels. All four are console-visible and archived in JSONL.}
\end{table}

\subsection{The KernelEventLog}

The \texttt{KernelEventLog} is the runtime substrate of the continual printing
system.  It is the single source of truth for all downstream analysis,
reporting, and dashboard records.

\begin{principle}
  The \texttt{KernelEventLog} owns the simulation history.  The
  \texttt{.xyz}/\texttt{.xyzFull} files own the atomistic ground truth.
  No other data structure may claim to represent either.  These two
  concerns do not overlap.
\end{principle}

The event log feeds the downstream pipeline in a single directed chain:

\[
\texttt{KernelEventLog}
  \;\to\;
  \texttt{FormationRecord}
  \;\to\;
  \texttt{run\_pipeline()}
  \;\to\;
  \texttt{AnalysisRecord}
  \;\to\;
  \texttt{ReportRecord}
  \;\to\;
  \texttt{DashboardRecord}
\]

The IKK analysis layer taps into this chain at the \texttt{AnalysisRecord}
boundary:

\[
\texttt{AnalysisRecord}
  \;\xrightarrow{\;\text{WO-75A}\;}
  \texttt{IKKEndTag}
  \;\xrightarrow{\;\text{render}\;}
  \texttt{ikk\_end\_tag\_md},\;\texttt{ikk\_end\_tag\_tex}
\]

% ===========================================================================
\section{IKK End-Tag Enrichment (WO-75A, Part~A)}
% ===========================================================================

\subsection{Motivation}

Before WO-75A, the simulation's identity metrics---projection loss,
hidden residual, distinguishability fraction---were computed and stored in
\texttt{IdentitySidecarRecord} but never \emph{explained}.
Reports read as data dumps.
The IKK End-Tag closes every report section with a structured block that
interprets these numbers in the language of IKK theory
(IKK I--IV, \textsc{Notation Registry v1.0}).

\subsection{The End-Tag Block Structure}

Each end-tag block emits the following fields (Markdown and \LaTeX{} variants):

\begin{table}[H]
\centering
\small
\begin{tabular}{@{}lll@{}}
\toprule
Field & Symbol & Meaning \\
\midrule
Scale regime      & $\Sn$          & Atomistic / molecular / material regime \\
$D$ mean          & $\Drec$        & Recoverable distinguishability fraction $\in[0,1]$ \\
$D$ loss          & $\DeltaD$      & Identity content lost across this run \\
Projection loss   &                & Loss from atomistic $\to$ observable projection \\
Hidden residual   & $\Psihid$      & Non-recoverable identity content estimate \\
Efficiency        & $\etaab$       & Cross-scale recovery fraction $\in[0,1]$ \\
Entropy delta     & $\DeltaS$      & Proportional to $-\DeltaD$ (IKK 2nd law) \\
Formation memory  &                & Fraction of formation history still recoverable \\
Second-law check  &                & \Dpass{PASS} / \Dwarn{WARN} / \Dfail{FAIL} \\
\bottomrule
\end{tabular}
\caption{IKK End-Tag fields. All values are derived from \texttt{IdentitySidecarRecord}---never
  written back into particle state or \texttt{.xyz} truth files.}
\end{table}

\subsection{Second-Law Check}

The second-law badge is a mandatory assertion appended to every section:

\begin{definition}[IKK Second Law]\label{def:secondlaw}
  $\DeltaS > 0$ if and only if $\DeltaD < 0$.
  Identity loss is monotonically paired with entropy increase in a
  non-reversible projection across scales.
\end{definition}

Badge thresholds (configurable via \texttt{[analysis.ikk\_end\_tag]}):

\begin{itemize}
  \item \Dpass{PASS}: $\Drec \geq 0.60$ and second-law consistent.
  \item \Dwarn{WARN}: $0.35 \leq \Drec < 0.60$ or mild second-law tension.
  \item \Dfail{FAIL}: $\Drec < 0.35$ or second-law violated.
\end{itemize}

\subsection{Implementation Summary (Day 75)}

\begin{table}[H]
\centering\small
\begin{tabular}{@{}lll@{}}
\toprule
Deliverable & File & Status \\
\midrule
\texttt{VsimIkkEndTagSection} struct
  & \texttt{include/vsim/vsim\_document.hpp}
  & \textcolor{vsimgreen}{\textbf{DONE}} \\
\texttt{apply\_ikk\_end\_tag\_key()} parser wiring
  & \texttt{src/vsim/vsim\_parser.cpp}
  & \textcolor{vsimgreen}{\textbf{DONE}} \\
\texttt{IKKEndTag} struct + render functions
  & \texttt{include/vsim/analysis/ikk\_end\_tag.hpp}
  & \textcolor{vsimgreen}{\textbf{DONE}} \\
\texttt{IkkEndTagModule} self-registering class
  & \texttt{src/vsim/analysis/ikk\_end\_tag.cpp}
  & \textcolor{vsimgreen}{\textbf{DONE}} \\
\texttt{AnalysisRecord} output fields
  & \texttt{include/vsim/analysis/i\_analysis\_module.hpp}
  & \textcolor{vsimgreen}{\textbf{DONE}} \\
Group~87 tests (20/20 PASS)
  & \texttt{tests/test\_ikk\_end\_tag.cpp}
  & \textcolor{vsimgreen}{\textbf{DONE}} \\
\midrule
GL D-colour overlay (Part~B, \S\ref{sec:partb})
  & \texttt{src/vis/renderer.cpp}
  & \textcolor{vsimamber}{\textbf{DEFERRED}} \\
\bottomrule
\end{tabular}
\caption{WO-75A Part~A deliverables. Part~B is fully specified but deferred to the next arc.}
\end{table}

\subsection{VSIM Script Surface}

\begin{lstlisting}[style=eventprint,
  caption={\texttt{[analysis.ikk\_end\_tag]} block in a \texttt{.vsim} script.}]
[analysis.ikk_end_tag]
emit_markdown      = true
emit_latex         = false
d_pass_threshold   = 0.60
d_warn_threshold   = 0.35
scale_regime       = "S3"
section_reference  = "IKK IV sec 3.2"
\end{lstlisting}

% ===========================================================================
\section{WO-75A Part~B --- GL D-Colour Overlay (Deferred)}\label{sec:partb}
% ===========================================================================

\begin{mdframed}[style=warnbox]
\textbf{Deferred to next arc.}\quad
The specification below is complete and frozen.  Nothing has been removed
from the design; only execution is deferred.  The next arc begins here.
\end{mdframed}

\subsection{Colour Mapping Convention}

The overlay maps $\Drec \in [0,1]$ to a cold-to-hot palette following IKK
convention (cold = identity preserved; hot = identity dissolved):

\[
\Drec = 1.0 \;\longrightarrow\; \text{deep blue}
\qquad
\Drec = 0.75 \;\longrightarrow\; \text{cyan}
\qquad
\Drec = 0.5 \;\longrightarrow\; \text{green}
\]
\[
\Drec = 0.25 \;\longrightarrow\; \text{amber}
\qquad
\Drec = 0.0 \;\longrightarrow\; \text{red}
\]

\subsection{Renderer Surface (Specified)}

\begin{lstlisting}[style=eventprint,
  caption={New renderer API surface -- WO-75A Part~B.}]
enum class RenderPass {
	STANDARD, DIST_OVERLAY, HIDDEN_RESIDUAL, SCALE_LADDER
};

void set_render_pass(RenderPass pass);
void set_dist_data(const std::vector<float>& per_atom_dist);
void set_hidden_residual_data(const std::vector<float>& per_atom_hid);
void render_dist_legend(int width, int height);
void render_scale_ladder_bar(int width, int height,
							 const std::string& active_regime);
\end{lstlisting}

Toggle: \texttt{D} key cycles
$\texttt{STANDARD} \to \texttt{DIST\_OVERLAY}
 \to \texttt{HIDDEN\_RESIDUAL} \to \texttt{SCALE\_LADDER}
 \to \texttt{STANDARD}$.

\subsection{Scale Ladder Panel}

\begin{lstlisting}[style=eventprint,
  caption={Scale ladder bottom bar concept.}]
S0(existence) --- S2(colour/EM) --- S3(spatial/EM) --- S4(weak) --- S_mat
										   ^ YOU ARE HERE
	eta_ab = 0.91   |Psi^hid| = 0.13   Delta-D = -0.11
\end{lstlisting}

% ===========================================================================
\section{IKK Identity Vector (WO-75B)}
% ===========================================================================

\subsection{Overview}

The IKK Identity Vector $\ivec$ is a five-component frame record that maps
the simulation state at frame $f$ into the IKK identity-information space:

\[
\ivec
=
\begin{pmatrix}
  x \\ y \\ z \\ t \\ w
\end{pmatrix}
=
\begin{pmatrix}
  \text{existence (presence)} \\
  \text{EM / structural signature} \\
  \text{spatial distinguishability} \\
  \text{temporal coherence} \\
  \text{internal / hidden channel}
\end{pmatrix}
\]

The vector is a \emph{proxy} mapping in Phase~1 (v5.13.5); the full
field-theory mapping is targeted for a future arc.

\subsection{IKKIdentitySeries}

Each run produces an \texttt{IKKIdentitySeries}: an ordered list of
\texttt{IKKIdentityFrameRecord} entries, one per simulation frame.
Each record carries $\ivec$, the frame-level $\DeltaD$, and an entropy proxy.

\begin{table}[H]
\centering\small
\begin{tabular}{@{}lll@{}}
\toprule
Deliverable & File & Status \\
\midrule
\texttt{IKKIdentityVector} + \texttt{IKKIdentitySeries}
  & \texttt{include/vsim/analysis/ikk\_identity\_vector.hpp}
  & \textcolor{vsimgreen}{\textbf{DONE}} \\
\texttt{from\_sidecar\_record()} Phase~1 proxy mapping
  & same
  & \textcolor{vsimgreen}{\textbf{DONE}} \\
\texttt{build\_ivec\_series()} + \texttt{write\_identity\_json()}
  & \texttt{src/vsim/analysis/ikk\_identity\_vector.cpp}
  & \textcolor{vsimgreen}{\textbf{DONE}} \\
Group~88 tests (28/28 PASS)
  & \texttt{tests/test\_ikk\_identity\_vector.cpp}
  & \textcolor{vsimgreen}{\textbf{DONE}} \\
\midrule
Phase~2 field-theory mapping
  & TBD
  & \textcolor{vsimamber}{\textbf{NEXT ARC}} \\
\bottomrule
\end{tabular}
\caption{WO-75B deliverables at the v5.13 freeze.}
\end{table}

\subsection{VSIM Script Surface}

\begin{lstlisting}[style=eventprint,
  caption={\texttt{[analysis.ivec]} block in a \texttt{.vsim} script.}]
[analysis.ivec]
enabled             = true
write_identity_json = true
compute_ivec_flux   = false   # J_iv -- identity-vector crossing signature
\end{lstlisting}

% ===========================================================================
\section{D Time-Series Plot (WO-75A-C) --- Delivered}
% ===========================================================================

\texttt{plot\_dist\_timeseries()} was added to
\texttt{reporting/generate\_report.py} and produces a dual-axis matplotlib
figure embedded in every consolidated report:

\begin{itemize}
  \item \textbf{Left axis:} $\Drec$ over simulation steps (blue line).
  \item \textbf{Right axis:} entropy proxy $\DeltaS$ (red dashed line).
  \item \textbf{Amber band:} region where $\DeltaD < 0$ (identity loss zone).
  \item \textbf{Reference line:} $D = 0.5$ (IKK~IV property activation threshold).
\end{itemize}

Output: \texttt{out/<run\_id>/dist\_timeseries.png}, embedded in
\texttt{report.md} with a caption auto-generated from the IKK end-tag fields.

% ===========================================================================
\section{v5.13 Branch Freeze Summary}
% ===========================================================================

\subsection{Test Registry at Freeze}

\begin{table}[H]
\centering\small
\begin{tabular}{@{}lrr@{}}
\toprule
Scope & Groups & Tests \\
\midrule
Core / Geometry / VSEPR (Groups 1--22)        & 22 & --- \\
Pipeline / Dashboard / Parser (Groups 23--43) & 21 & --- \\
EigenMine / CurveFit / Release (44--46)       &  3 &  10 \\
CTL pipeline (Groups 58--69)                  & 12 & 106 \\
Dynx / DUB / Demo pipeline (71--77)           &  7 & 113 \\
LengthScale / OutputFilter / GapClass (78--80)&  3 & --- \\
ObserveMetrics / FormationHelper / Bio (81--83)&  3 & --- \\
Replay / Discovery / CrystalMD (84--86)       &  3 & --- \\
IKK End-Tag (Group 87)                        &  1 &  20 \\
IKK Identity Vector (Group 88)                &  1 &  28 \\
MCF-CAI (Group 89)                            &  1 & --- \\
\midrule
\textbf{Total at freeze}
  & \textbf{89}
  & \textbf{167 / 167\quad PASS} \\
\bottomrule
\end{tabular}
\caption{Test registry at the v5.13 freeze (Day~82, commit \texttt{489d11cd}).}
\end{table}

\subsection{Arc Delivery Table}

\begin{table}[H]
\centering\small
\begin{tabular}{@{}lll@{}}
\toprule
Version & Day & Deliverable \\
\midrule
v5.0.0-beta.7  & 57 & render\_interval / step emission cadence \\
v5.0.0-beta.8  & 58 & PBC / cell / boundary / Ewald \\
v5.0.0-beta.9  & 59 & Registry resolution engine / CLI layer \\
v5.0.0-beta.10 & 60 & Variance / N\_evolution / while / batch sweep \\
v5.0.0-beta.11 & 61 & Macro sampling / empirical verification \\
v5.0.0-beta.12 & 62 & Batching upgrade / study orchestration \\
v5.1.13        & 72 & Pillar F: empirical chemistry, Dynx, CTL pipeline \\
v5.13.3        & 73 & Wizard modules 6--8 full field coverage \\
v5.13.4        & 74 & Gap classifier + output filter + chemistry audit \\
v5.13.5        & 75 & IKK End-Tag (WO-75A-A), Identity Vector (WO-75B), MCF-CAI (WO-76) \\
\midrule
\textbf{v5.13 FROZEN} & \textbf{82}
  & \textbf{Branch frozen. Chapter~20 authored as freeze artefact.} \\
\bottomrule
\end{tabular}
\caption{v5.13 arc delivery table. Work commenced at Day~57 (v5.0.0-beta.7).}
\end{table}

\subsection{Deprecated at Freeze}

\begin{table}[H]
\centering\small
\begin{tabular}{@{}lll@{}}
\toprule
ID & Item & Rationale \\
\midrule
MF-F02
  & GitHub Actions CI workflow
  & Out of scope for v5.13. Not blocking any UC\@. Deferred indefinitely. \\
\bottomrule
\end{tabular}
\caption{Items deprecated at the v5.13 freeze.}
\end{table}

\subsection{Deferred to Next Arc}

\begin{table}[H]
\centering\small
\begin{tabular}{@{}lll@{}}
\toprule
WO & Item & Next-arc entry point \\
\midrule
WO-75A-B
  & IKK GL D-colour overlay
  & \texttt{src/vis/renderer.cpp}: add \texttt{RENDER\_PASS\_DIST} \\
WO-75B gate
  & \texttt{vsepr doctor} full pass + tag push
  & Runtime data path audit \\
MF-B01
  & \texttt{[sweep]} runtime dispatch
  & \texttt{src/vsim/vsim\_runtime.hpp}: sweep executor \\
\bottomrule
\end{tabular}
\caption{Items deferred to the next arc. All specifications are complete and frozen here.}
\end{table}

% ===========================================================================
\section*{Closing Note}
% ===========================================================================

\addcontentsline{toc}{section}{Closing Note}

The v5.13 arc set out to formalise a single idea that the project had been
building toward since Day~1: that a simulation should not merely
\emph{compute} but \emph{continuously account for what it knows and what it
has lost}.

This chapter documents that arc's architectural outcome.  The Continual
Printing System is now a named, documented doctrine.  The IKK layer has been
wired from theory (IKK I--IV) through schema
(\texttt{VsimIkkEndTagSection}, \texttt{[analysis.ivec]}) into the runtime
(\texttt{IkkEndTagModule}, \texttt{build\_ivec\_series()}), reporting
(\texttt{render\_ikk\_end\_tag\_md()}, \texttt{plot\_dist\_timeseries()}),
and test coverage (Groups 87--89, 48$+$ cases).

The GL overlay (Part~B) is the one remaining visual expression of this
system.  Its specification is complete.  The next arc begins there.

\vspace{16pt}
\begin{mdframed}[style=freezebox]
\centering
\textbf{v5.13 arc --- FROZEN}
\quad|\quad Day~82
\quad|\quad 167/167 tests
\quad|\quad Commit \texttt{489d11cd}\\[4pt]
\small
Next arc opens from this state.\quad
First task: \textbf{WO-75A-B} ---
wire \texttt{RENDER\_PASS\_DIST} into \texttt{src/vis/renderer.cpp}.
\end{mdframed}

\end{document}
